% GENERATED FILE. Edit canonical lessons, then run npm run generate:books.
% canonical-source-hash: fnv1a64:2ead0adb62f332cb
% canonical-lessons: GU-C506-vankar, GU-C506-tatno, GU-C506-vanat, GU-C506-tanvu, GU-C506-vintvu

\chapter{At the Loom}
\label{ch:gu-at-the-loom}
\hlchaptermodality{506}

\begin{chapteropening}
\textbf{By the end of this chapter:} \emph{I can talk about weaving cloth on a loom.}

Complete the last lesson of chapter 506: I can talk about weaving cloth on a loom.
\end{chapteropening}

\section[\texorpdfstring{\gu{વણકર} (vaṇkar)}{વણકર (vaṇkar)}]{\gu{વણકર} (vaṇkar) — a weaver}

\label{lesson:GU-C506-vankar}



\begin{quote}
Before the new one: say the Gujarati for to spin (thread), then the Gujarati for khadi.
\end{quote}

\subsection*{You'll want to know: \gu{વણકર}}
\textbf{\gu{વણકર}} — \emph{vaṇkar} — \textquotedblleft{}a weaver\textquotedblright{}.

\begin{quote}
\gu{આ} \gu{ગામના} \gu{વણકર} \gu{સુંદર} \gu{સાડીઓ} \gu{બનાવે} \gu{છે}.
\end{quote}

\begin{quote}
\emph{ā gāmnā vaṇkar sundar sāḍīo banāve che.}
\end{quote}

\begin{quote}
The weavers of this village make beautiful saris.
\end{quote}

A masculine or feminine noun, by the person meant.

\begin{cousinweb}
\textbf{\gu{વણકર}} (\emph{vaṇkar}) is built on \emph{vaṇvũ}, \textquotedblleft{}to weave\textquotedblright{}, which CDIAL lists under \emph{*vināti} (CDIAL 11773), with the ending \emph{-kar}, \textquotedblleft{}one who does\textquotedblright{}.
\end{cousinweb}

\begin{grammarlens}[title={What you've built}]
One more word for describing what you lived and felt.
\end{grammarlens}

\subsection*{Your turn}
\begin{itemize}
\raggedright
  \item \emph{Say it:} \emph{vaṇkar}
  \item \emph{Say it:} \emph{vaṇkar}, once more
  \item \emph{Say it:} \emph{ā gāmnā vaṇkar sundar sāḍīo banāve che.} — The weavers of this village make beautiful saris.
  \item \emph{Recall:} say the Gujarati for to spin (thread), then the Gujarati for khadi, then say \emph{vaṇkar} again
\end{itemize}

\begin{tcolorbox}[breakable,colback=teal!4,colframe=teal!35!black,title={Before you move on}]
What does \gu{વણકર} mean? (\textquotedblleft{}A weaver\textquotedblright{}.) Say it once more.
\end{tcolorbox}
\section[\texorpdfstring{\gu{તાંતણો} (tā̃tṇo)}{તાંતણો (tā̃tṇo)}]{\gu{તાંતણો} (tā̃tṇo) — a strand (of fibre)}

\label{lesson:GU-C506-tatno}



\begin{quote}
Before the new one: say the Gujarati for khadi, then the Gujarati for a weaver.
\end{quote}

\subsection*{You'll want to know: \gu{તાંતણો}}
\textbf{\gu{તાંતણો}} — \emph{tā̃tṇo} — \textquotedblleft{}a strand (of fibre)\textquotedblright{}.

\begin{quote}
\gu{રૂના} \gu{તાંતણામાંથી} \gu{દોરો} \gu{બને} \gu{છે}.
\end{quote}

\begin{quote}
\emph{rūnā tā̃tṇāmā̃thī doro bane che.}
\end{quote}

\begin{quote}
Thread is made from strands of cotton.
\end{quote}

A masculine noun.

\begin{cousinweb}
\textbf{\gu{તાંતણો}} (\emph{tā̃tṇo}) is inherited from Sanskrit \emph{tántu}, \textquotedblleft{}a thread, the warp\textquotedblright{}, and CDIAL lists the Gujarati word with the sense \textquotedblleft{}thread\textquotedblright{} (CDIAL 5661). Before a postposition it becomes \textbf{\gu{તાંતણા}-}.
\end{cousinweb}

\begin{grammarlens}[title={What you've built}]
One more word for describing what you lived and felt.
\end{grammarlens}

\subsection*{Your turn}
\begin{itemize}
\raggedright
  \item \emph{Say it:} \emph{tā̃tṇo}
  \item \emph{Say it:} \emph{tā̃tṇo}, once more
  \item \emph{Say it:} \emph{rūnā tā̃tṇāmā̃thī doro bane che.} — Thread is made from strands of cotton.
  \item \emph{Recall:} say the Gujarati for khadi, then the Gujarati for a weaver, then say \emph{tā̃tṇo} again
\end{itemize}

\begin{tcolorbox}[breakable,colback=teal!4,colframe=teal!35!black,title={Before you move on}]
What does \gu{તાંતણો} mean? (\textquotedblleft{}A strand (of fibre)\textquotedblright{}.) Say it once more.
\end{tcolorbox}
\section[\texorpdfstring{\gu{વણાટ} (vaṇāṭ)}{વણાટ (vaṇāṭ)}]{\gu{વણાટ} (vaṇāṭ) — a weave}

\label{lesson:GU-C506-vanat}



\begin{quote}
Before the new one: say the Gujarati for a weaver, then the Gujarati for a strand.
\end{quote}

\subsection*{You'll want to know: \gu{વણાટ}}
\textbf{\gu{વણાટ}} — \emph{vaṇāṭ} — \textquotedblleft{}a weave\textquotedblright{}.

\begin{quote}
\gu{આ} \gu{સાડીનો} \gu{વણાટ} \gu{ઘણો} \gu{સુંદર} \gu{છે}.
\end{quote}

\begin{quote}
\emph{ā sāḍīno vaṇāṭ ghaṇo sundar che.}
\end{quote}

\begin{quote}
The weave of this sari is very beautiful.
\end{quote}

A masculine noun.

\begin{cousinweb}
\textbf{\gu{વણાટ}} (\emph{vaṇāṭ}) is built on \emph{vaṇvũ}, \textquotedblleft{}to weave\textquotedblright{}, which CDIAL lists under \emph{*vināti} (CDIAL 11773), with the ending \emph{-āṭ}: the way a cloth is woven.
\end{cousinweb}

\begin{grammarlens}[title={What you've built}]
One more word for describing what you lived and felt.
\end{grammarlens}

\subsection*{Your turn}
\begin{itemize}
\raggedright
  \item \emph{Say it:} \emph{vaṇāṭ}
  \item \emph{Say it:} \emph{vaṇāṭ}, once more
  \item \emph{Say it:} \emph{ā sāḍīno vaṇāṭ ghaṇo sundar che.} — The weave of this sari is very beautiful.
  \item \emph{Recall:} say the Gujarati for a weaver, then the Gujarati for a strand, then say \emph{vaṇāṭ} again
\end{itemize}

\begin{tcolorbox}[breakable,colback=teal!4,colframe=teal!35!black,title={Before you move on}]
What does \gu{વણાટ} mean? (\textquotedblleft{}A weave\textquotedblright{}.) Say it once more.
\end{tcolorbox}
\section[\texorpdfstring{\gu{તાણવું} (tāṇvũ)}{તાણવું (tāṇvũ)}]{\gu{તાણવું} (tāṇvũ) — to stretch tight}

\label{lesson:GU-C506-tanvu}



\begin{quote}
Before the new one: say the Gujarati for a strand, then the Gujarati for a weave.
\end{quote}

\subsection*{You'll want to know: \gu{તાણવું}}
\textbf{\gu{તાણવું}} — \emph{tāṇvũ} — \textquotedblleft{}to stretch tight\textquotedblright{}.

\begin{quote}
\gu{વણકરે} \gu{બે} \gu{લાકડાં} \gu{વચ્ચે} \gu{દોરા} \gu{તાણ્યા}.
\end{quote}

\begin{quote}
\emph{vaṇakre be lākḍā̃ vacce dorā tāṇyā.}
\end{quote}

\begin{quote}
The weaver stretched the threads tight between two poles.
\end{quote}

Past, with \gu{તેણે} (he or she …): \gu{તેણે} \gu{દોરો} \gu{તાણ્યો}, \emph{teṇe doro tāṇyo}.

\begin{cousinweb}
\textbf{\gu{તાણવું}} (\emph{tāṇvũ}) is inherited from Sanskrit \emph{tānáyati}, \textquotedblleft{}makes stretch\textquotedblright{}, and CDIAL lists the Gujarati verb, \textquotedblleft{}to stretch, pull\textquotedblright{} (CDIAL 5762).
\end{cousinweb}

\begin{grammarlens}[title={What you've built}]
One more word for describing what you lived and felt.
\end{grammarlens}

\subsection*{Your turn}
\begin{itemize}
\raggedright
  \item \emph{Say it:} \emph{tāṇvũ}
  \item \emph{Say it:} \emph{tāṇvũ}, once more
  \item \emph{Say it:} \emph{vaṇakre be lākḍā̃ vacce dorā tāṇyā.} — The weaver stretched the threads tight between two poles.
  \item \emph{Recall:} say the Gujarati for a strand, then the Gujarati for a weave, then say \emph{tāṇvũ} again
\end{itemize}

\begin{tcolorbox}[breakable,colback=teal!4,colframe=teal!35!black,title={Before you move on}]
What does \gu{તાણવું} mean? (\textquotedblleft{}To stretch tight\textquotedblright{}.) Say it once more.
\end{tcolorbox}
\section[\texorpdfstring{\gu{વીંટવું} (vī̃ṭvũ)}{વીંટવું (vī̃ṭvũ)}]{\gu{વીંટવું} (vī̃ṭvũ) — to wind round (a thread)}

\label{lesson:GU-C506-vintvu}



\begin{quote}
Before the new one: say the Gujarati for a weave, then the Gujarati for to stretch tight.
\end{quote}

\subsection*{You'll want to know: \gu{વીંટવું}}
\textbf{\gu{વીંટવું}} — \emph{vī̃ṭvũ} — \textquotedblleft{}to wind round (a thread)\textquotedblright{}.

\begin{quote}
\gu{મીરાએ} \gu{દોરો} \gu{આંગળી} \gu{પર} \gu{વીંટ્યો}.
\end{quote}

\begin{quote}
\emph{Mīrāe doro ā̃gaḷī par vī̃ṭyo.}
\end{quote}

\begin{quote}
Mira wound the thread round her finger.
\end{quote}

Past, with \gu{તેણે} (he or she …): \gu{તેણે} \gu{દોરો} \gu{વીંટ્યો}, \emph{teṇe doro vī̃ṭyo}.

\begin{cousinweb}
\textbf{\gu{વીંટવું}} (\emph{vī̃ṭvũ}) is filed by CDIAL under Sanskrit \emph{vīṭā}, a small rolled piece, and CDIAL lists the Gujarati verb, \textquotedblleft{}to wrap, roll, wind\textquotedblright{} (CDIAL 12045). The ring \textbf{\gu{વીંટી}} belongs to the same family.
\end{cousinweb}

\begin{grammarlens}[title={What you've built}]
One more word for describing what you lived and felt.
\end{grammarlens}

\subsection*{Your turn}
\begin{itemize}
\raggedright
  \item \emph{Say it:} \emph{vī̃ṭvũ}
  \item \emph{Say it:} \emph{vī̃ṭvũ}, once more
  \item \emph{Say it:} \emph{Mīrāe doro ā̃gaḷī par vī̃ṭyo.} — Mira wound the thread round her finger.
  \item \emph{Recall:} say the Gujarati for a weave, then the Gujarati for to stretch tight, then say \emph{vī̃ṭvũ} again
\end{itemize}

\begin{tcolorbox}[breakable,colback=teal!4,colframe=teal!35!black,title={Before you move on}]
What does \gu{વીંટવું} mean? (\textquotedblleft{}To wind round (a thread)\textquotedblright{}.) Say it once more.
\end{tcolorbox}
